% GENERATED FILE. Edit canonical lessons, then run npm run generate:books.
% canonical-source-hash: fnv1a64:15444533f20d5cd7
% canonical-lessons: FR-C196-boucher, FR-C196-mecanicien, FR-C196-electricien, FR-C196-plombier, FR-C196-ouvrier

\chapter{Butcher, Mechanic, Electrician, Plumber, Worker}
\label{ch:fr-butcher-mechanic-electrician-plumber-worker}
\hlchaptermodality{196}

\begin{chapteropening}
\textbf{By the end of this chapter:} \emph{I can say a butcher, a mechanic, an electrician, a plumber and a worker.}

Complete the last lesson of chapter 196: I can say a butcher, a mechanic, an electrician, a plumber and a worker.
\end{chapteropening}

\section[le boucher]{le boucher — a butcher}

\label{lesson:FR-C196-boucher}



\begin{quote}
Before the new one: say the French for a screen, then the French for a keyboard.
\end{quote}

\subsection*{You'll want to know: le boucher}
\textbf{le boucher} — \textquotedblleft{}a butcher\textquotedblright{}.

\begin{grammarlens}[title={What you've built}]
One more word for this chapter.
\end{grammarlens}

\subsection*{Your turn}
\begin{itemize}
\raggedright
  \item \emph{Say it:} \emph{le boucher}
  \item \emph{Say it:} \emph{le boucher}, once more
  \item \emph{Say it:} say \emph{le boucher}
  \item \emph{Recall:} say the French for a screen, then the French for a keyboard, then say \emph{le boucher} again
\end{itemize}

\begin{tcolorbox}[breakable,colback=teal!4,colframe=teal!35!black,title={Before you move on}]
What does le boucher mean? (\textquotedblleft{}A butcher\textquotedblright{}.) Say it once more.
\end{tcolorbox}
\section[le mécanicien]{le mécanicien — a mechanic}

\label{lesson:FR-C196-mecanicien}



\begin{quote}
Before the new one: say the French for a keyboard, then the French for a butcher.
\end{quote}

\subsection*{You'll want to know: le mécanicien}
\textbf{le mécanicien} — \textquotedblleft{}a mechanic\textquotedblright{}.

\begin{grammarlens}[title={What you've built}]
One more word for this chapter.
\end{grammarlens}

\subsection*{Your turn}
\begin{itemize}
\raggedright
  \item \emph{Say it:} \emph{le mécanicien}
  \item \emph{Say it:} \emph{le mécanicien}, once more
  \item \emph{Say it:} say \emph{le mécanicien}
  \item \emph{Recall:} say the French for a keyboard, then the French for a butcher, then say \emph{le mécanicien} again
\end{itemize}

\begin{tcolorbox}[breakable,colback=teal!4,colframe=teal!35!black,title={Before you move on}]
What does le mécanicien mean? (\textquotedblleft{}A mechanic\textquotedblright{}.) Say it once more.
\end{tcolorbox}
\section[l'électricien]{l'électricien — an electrician}

\label{lesson:FR-C196-electricien}



\begin{quote}
Before the new one: say the French for a butcher, then the French for a mechanic.
\end{quote}

\subsection*{You'll want to know: l'électricien}
\textbf{l'électricien} — \textquotedblleft{}an electrician\textquotedblright{}.

\begin{grammarlens}[title={What you've built}]
One more word for this chapter.
\end{grammarlens}

\subsection*{Your turn}
\begin{itemize}
\raggedright
  \item \emph{Say it:} \emph{l'électricien}
  \item \emph{Say it:} \emph{l'électricien}, once more
  \item \emph{Say it:} say \emph{l'électricien}
  \item \emph{Recall:} say the French for a butcher, then the French for a mechanic, then say \emph{l'électricien} again
\end{itemize}

\begin{tcolorbox}[breakable,colback=teal!4,colframe=teal!35!black,title={Before you move on}]
What does l'électricien mean? (\textquotedblleft{}An electrician\textquotedblright{}.) Say it once more.
\end{tcolorbox}
\section[le plombier]{le plombier — a plumber}

\label{lesson:FR-C196-plombier}



\begin{quote}
Before the new one: say the French for a mechanic, then the French for an electrician.
\end{quote}

\subsection*{You'll want to know: le plombier}
\textbf{le plombier} — \textquotedblleft{}a plumber\textquotedblright{}.

\begin{grammarlens}[title={What you've built}]
One more word for this chapter.
\end{grammarlens}

\subsection*{Your turn}
\begin{itemize}
\raggedright
  \item \emph{Say it:} \emph{le plombier}
  \item \emph{Say it:} \emph{le plombier}, once more
  \item \emph{Say it:} say \emph{le plombier}
  \item \emph{Recall:} say the French for a mechanic, then the French for an electrician, then say \emph{le plombier} again
\end{itemize}

\begin{tcolorbox}[breakable,colback=teal!4,colframe=teal!35!black,title={Before you move on}]
What does le plombier mean? (\textquotedblleft{}A plumber\textquotedblright{}.) Say it once more.
\end{tcolorbox}
\section[l'ouvrier]{l'ouvrier — a worker}

\label{lesson:FR-C196-ouvrier}



\begin{quote}
Before the new one: say the French for an electrician, then the French for a plumber.
\end{quote}

\subsection*{You'll want to know: l'ouvrier}
\textbf{l'ouvrier} — \textquotedblleft{}a worker\textquotedblright{}.

\begin{grammarlens}[title={What you've built}]
One more word for this chapter.
\end{grammarlens}

\subsection*{Your turn}
\begin{itemize}
\raggedright
  \item \emph{Say it:} \emph{l'ouvrier}
  \item \emph{Say it:} \emph{l'ouvrier}, once more
  \item \emph{Say it:} say \emph{l'ouvrier}
  \item \emph{Recall:} say the French for an electrician, then the French for a plumber, then say \emph{l'ouvrier} again
\end{itemize}

\begin{tcolorbox}[breakable,colback=teal!4,colframe=teal!35!black,title={Before you move on}]
What does l'ouvrier mean? (\textquotedblleft{}A worker\textquotedblright{}.) Say it once more.
\end{tcolorbox}
